\chapter{Related Work\label{cha:chapter5}}

In this chapter, we survey work related to the thesis's contributions in four directions: 

\begin{itemize}
    \item hardware-conscious index structures, 
    \item vectorised and parallel query processing, 
    \item approximate histogram-based systems, and 
    \item the performance engineering of Python-based data systems.
\end{itemize}

\section{Hardware-Optimised Index Structures}

Work done by Manegold et al.~\cite{manegold_2000} illustrates that main-memory access is the dominant bottleneck for in-memory database operations, from that work cache-conscious design principles were established, which subsequently influenced much of the work in database systems. Manegold et al.'s analysis showed that column-major (SoA) storage dramatically outperforms row-major (AoS) storage for projective (columnar) access patterns, which is the direct motivation for the memory layout designed in Chapter~\ref{cha:chapter3}.

Adaptive Radix Trees (ART)~\cite{leis_2013} adapt their node fanout to the key distribution, achieving cache efficiency for point and range queries. They avoid fixed-fanout waste of classical B-trees. ARTs target key-value loops rather than percentile predicate evaluations, nonetheless their design philosophy of having a data structure based on cache line behaviour rather than asymptotic complexity informs the approach taken in this thesis.

Khuong and Morin~\cite{khuong_morin_2017} established in a survey for comparison-based searching on modern CPUs that the Eytzinger layout (which is breadth-first) outperforms classical sorted arrays when workloads are cache-cold and of moderate size, due to its cache-line-aligned root of the search tree. This analysis was extended by Schlegel et al.~\cite{schlegel_2009} with SIMD-accelerated k-ary search. Neither of these constructions survived into the final integrated Fainder build, due to Fainder's elimination of cache-coldness on L1-resident columns and the k-ary chain shortening landing on the slack outside of $t=1$. The foundational analysis behind these designs traces back to Brodal and Fagerberg~\cite{brodal_fagerberg_2006}.

A learned index structure~\cite{kraska_2018} replaces B-tree index nodes with learned models that predict key positions. This trades worst-case guarantees (e.g. logarithmic search time) for expected-case cache efficiency. The problem with learned indexes is that they are only effective on learnable key distributions, that is documented by Kraska et al.~\cite{kraska_2018}. Classical sorted-array search with layouts that are hardware-conscious outperform these types of indexes when the key distribution is heterogeneous and cannot be accurately approximated by a small model. Fainder's histogram collection is exactly a heterogeneous workload, which is why classical binary search with SoA layout is the right choice.


\section{Vectorised and Parallel Query Processing}

For analytical queries, Neumann~\cite{neumann_2011} demonstrated that compiling query plans to native machine code rather than interpreting them through a tuple-at-a-time Volcano model yields order-of-magnitude performance improvements. Key insight here is that inner loops that are tight and specialised by type can be optimised by LLVM better than interpreted bytecode. The entire idea of Fainder's Rust implementation is motivated by this insight.

Kersten et al.~\cite{kersten_2018} compared compiled and vectorised query-execution models and found that neither dominates: vectorised execution wins where cache-resident batch processing and parallel data access matter, compilation where per-tuple work is tight. We implement explicit AVX-512 SIMD intrinsics for the final stage of Fainder's binary search (Section~\ref{sec:exp:c1:simd}), replacing the last ${\sim}16$ scalar iterations with a single vector comparison.

Leis et al.~\cite{leis_2014} showed with morsel-driven parallelism that dynamic work distribution outperforms static task partitioning when sub-tasks have unequal cost. Fainder's cluster queries have costs that are heterogeneous in exactly this way (large cluster $\to$ more histograms $\to$ more binary search steps). This leads us to choose Rayon's work stealing scheduler (an implementation of the scheduler of Blumofe and Leiserson~\cite{blumofe_leiserson_1999}) over round-robin thread assignment. Our reports in Chapter~\ref{cha:chapter4} are consistent with this: speedup grows with thread count until the memory-traffic and interconnect ceilings bind (§\ref{sec:exp:scaling}).

Kocberber et al.~\cite{kocberber_2013} propose a way to hide memory latency in in-memory hash-join traversal. This is possible by interleaving many independent lookups per core, saturating the MSHR budget. Fainder's 8-way batch binary search ablation (Section~\ref{sec:exp:c1:batch}) is a direct application of this idea. 


The MonetDB/X100 project~\cite{boncz_2005} introduced vector-at-a-time execution, processing cache-resident batches to combine columnar layouts with tight SIMD-friendly inner loops.  Section~\ref{sec:exp:wf:columnar} is inspired by this, and the  design is surveyed in a broader sense by Abadi et al.~\cite{abadi_2008,abadi_2013}. Vectorised database operators studied by Willhalm et al.~\cite{willhalm_2009} and Zhou and Ross~\cite{zhou_ross_2002} were also influential in Fainder's new SIMD design.

\section{Approximate Histogram-Based Systems}

Distribution queries are often approximated through compact histogram sketches. Data structures such as T-digest~\cite{dunning_2019} and HdrHistogram~\cite{hdrhistogram} maintain compressed representations of \emph{streaming} data distributions that support approximate quantile queries. Synopsis-based query processing~\cite{ioannidis_2003,cormode_2011} precomputes compact statistical summaries to answer approximate aggregate queries without scanning raw data.

We do not consider these systems as competitors to Fainder, since Fainder uses precomputed histograms to answer \emph{exact} percentile predicates, whereas the above systems target \emph{approximate} quantile computation over streaming or continuously updated data. A better classification would be that these are complementary systems: Fainder could be used to answer exact percentile queries on the histograms produced by these systems, or these systems could be used to generate the histograms that Fainder indexes.

\section{Hardware-Aware Indexing of Adjacent Primitives}

Percentile-predicate evaluation is a finer grained primitive than the hardware-conscious systems surveyed above, and it has not received the same level of attention. Below are several closely related primitives, that help us position Fainder's percentile-predicate evaluation in the landscape of hardware-conscious systems.

\textbf{Cache-resident sorted-array search.} Fainder's binary search runs against cache-resident data after warm-up: the bin-edge array (${\sim}142$ floats per cluster) is L1-resident, and the per-cluster cumulative-density slices range from L1- to L2-resident depending on cluster density (§\ref{sec:method:data}). The literature on cache-conscious sorted-array search is as follows:

\begin{itemize}
    \item Kim et al.'s FAST index~\cite{kim_2010} reports large speedups by re-organising search trees to exploit page size, cache-line size, and SIMD width simultaneously.
    \item Khuong and Morin~\cite{khuong_morin_2017}'s work shows small $n$ branchless binary search wins, while Eytzinger BFS layout wins for large $n$.
    \item Schlegel et al.'s k-ary search on modern processors~\cite{schlegel_2009} is the foundation for the SIMD k-ary variants explored in earlier ablation campaigns of this thesis, and for the SIMD leaf stage of Fainder's binary search (Section~\ref{sec:exp:c1:simd}).
\end{itemize}

\textbf{Vectorised column scan.} Fainder's f16 ablation (Section~\ref{sec:exp:c2:f16}) is a vectorised scan of a column of 16-bit floats, which is the same access pattern as the following systems:
\begin{itemize}
    \item BitWeaving~\cite{li_patel_2013} achieves below-one-cycle-per-element scan throughput by packing column values across word boundaries.
    \item ByteSlice~\cite{feng_2015} refines this with a byte-aligned layout.
\end{itemize}

Wallewein-Eising et al.~\cite{wallewein_2018} survey SIMD acceleration across main-memory index structures more broadly.

\section{Composability of Optimisations in Database and Systems Engineering}

The negative-composability pattern of §\ref{sec:exp:negcomp} reasons with three performance engineering perspectives.

\paragraph{Cost-based and adaptive query optimisation.} Selinger et al.~\cite{selinger_1979} formulated query optimisation with costs that decompose additively across operators. Adaptive query processing~\cite{avnur_hellerstein_2000} keeps the additive model, but defers plan selection to runtime statistics. We refine the scope of these assumptions with our five negative composition instances (§\ref{sec:exp:negcomp:instances}): at a ceiling, the cost of optimisation $B$ depends on which ceiling optimisation $A$ has already collected, so we can't rely on per-operator costs. In this case, the pre-flight check plays the same structural role as the Selinger cost model, based on perf counters rather than cardinalities, we apply a falsifiable predictor before committing to implementation cost.

\paragraph{Hardware--software co-design.}

Caladan and heartbeat scheduling~\cite{fried_2020_caladan,acar_2018_heartbeat} manage hardware interference inside a decision loop. They share the view that hardware ceilings bind performance, instead of algorithmic primitives. However, their reasoning for composition differs: Caladan assumes that along orthogonal resources, interference is independently manageable, and heartbeat scheduling assumes that the work composes, and that coordination overhead can be bounded under a work--span analysis. Our instance of a cooperative-cache scheduler over an already-amortised baseline (§\ref{sec:exp:wf:morsel}) is a counter-example: it satisfies the work--span analysis and still regresses, because no residual L3 misses remain to recover.

\paragraph{The ``no silver bullet'' and memory-wall traditions.}
Gains are assumed to compose multiplicatively in the ``no silver bullet'' tradition of Brooks~\cite{brooks_1987} and Wulf and McKee~\cite{wulf_mckee_1995}. While Brooks argues that the axes have to be independent, Wulf and McKee argue that the axes are not independent, but that the only way to get past the memory wall is to apply multiple hardware-conscious techniques. We show yet again, with our five negative-composition instances, that the axes are not independent, but that the right optimisation at each composition step is contingent on which ceiling currently binds.

\paragraph{Positioning.}
Each instance in the composition study (§\ref{sec:exp:negcomp:instances}) had an a-priori argument under each of these three perspectives that the composition should be positive. Each instance was then falsified by a perf-counter-grounded mechanism. There is, to our knowledge, no prior documentation of this specific shape: five axes, mechanism-grounded ablation pairs, and a falsifiable pre-flight diagnostic.


\section{Performance Gap in Distribution-Aware Search}

The original Fainder paper~\cite{behme_2024} does not benchmark against a compiled implementation, nor does it investigate hardware-level efficiency. Instead, it benchmarks against profile-scan, NDIST, PSCAN, and BinSort baselines and focuses on algorithmic correctness and retrieval effectiveness. The SIGMOD demo~\cite{behme_2025_demo}, the companion theoretical framework work of Esmailpour et al.~\cite{esmailpour_2025}, and the data-discovery survey of Esmailoghli, Galhotra, and Abedjan~\cite{abedjan_2025} situate Fainder (or Fainder-like work) algorithmically without addressing systems performance. Related distribution-aware systems suffer the same drawback: Esmailoghli et al.~\cite{EsmailoghliQZ21} and Traub et al.~\cite{TraubGCBKRM20} propose alternative indexing strategies without analysing cache or parallelism behaviour, Goods~\cite{halevy_2016} treats retrieval at indexing scale, Zhang and Ives~\cite{zhang_2018} focus on algorithmic structure, and Asudeh et al.~\cite{asudeh_2022} motivate metadata-rich search from a fairness angle.

The contribution of this thesis addresses two gaps of this literature:

\begin{itemize}
    \item No studies, to our knowledge, have been conducted that characterise hardware bottlenecks of histogram-based percentile search or hardware-conscious engineering of such systems.
    \item Before this work, no evaluation existed, to our knowledge, of inter-query SIMD for binary-search-shaped workloads at Fainder's cluster-loop scale on Sapphire Rapids. Existing evaluations such as Polychroniou et al.\ and Lang et al.~\cite{polychroniou_2014,lang_2020} predate the microarchitecture.
\end{itemize}

This thesis addresses the first gap with Fainder's query pipeline, due to its embarrassingly parallel structure, its access pattern being binary search, and its read-only index. It closes the second directly: the horizontal-SIMD evaluation of §\ref{sec:exp:c1:horizontal} is, to our knowledge, the first such measurement on Sapphire Rapids, with a null result at cache-resident scale.